% Banque d'exercices — chapitre 8
% Le chapitre est déterminé par le nom de ce fichier.

\begin{exo}[Les indices ne sont pas optionnels]{derivees-partielles-indices}
\keywords{dérivées partielles, variables fixées, indices, isobare, isentropique, gaz parfait}

On veut savoir comment varie le volume de $N$ particules de gaz parfait monoatomique lorsque sa température augmente. La question est incomplète tant qu'on ne précise pas ce qui est maintenu constant.

\begin{enumerate}
\item Calculer $\left(\dfrac{\partial V}{\partial T}\right)_{P,N}$ (chauffage isobare).
\item Calculer $\left(\dfrac{\partial V}{\partial T}\right)_{S,N}$ (évolution isentropique). On pourra travailler en différentielles : que devient le premier principe quand $\mathrm{d}S = \mathrm{d}N = 0$ ?
\item Comparer les deux résultats (signes !) et interpréter physiquement.
\item Montrer que le rapport des deux dérivées vaut $-(\gamma - 1)$ et vérifier la cohérence avec la loi de Laplace $TV^{\gamma-1} = \text{cste}$.
\end{enumerate}

\begin{indication}
Question 1 : isoler $V$ dans $PV = Nk_BT$. Question 2 : $\mathrm{d}U = T\,\mathrm{d}S - P\,\mathrm{d}V = -P\,\mathrm{d}V$ et $\mathrm{d}U = \frac{3}{2}Nk_B\,\mathrm{d}T$.
\end{indication}

\begin{solution}
\textbf{Question 1.} À $P$ et $N$ fixés, $V = Nk_BT/P$ donc $\left(\dfrac{\partial V}{\partial T}\right)_{P,N} = \dfrac{Nk_B}{P} > 0$.

\textbf{Question 2.} À $S$ et $N$ fixés, le premier principe donne $\mathrm{d}U = -P\,\mathrm{d}V$, et par ailleurs $\mathrm{d}U = \tfrac{3}{2}Nk_B\,\mathrm{d}T$. Donc
\[
\left(\frac{\partial V}{\partial T}\right)_{S,N} = -\frac{3Nk_B}{2P} < 0.
\]

\textbf{Question 3.} La même question — « comment $V$ varie-t-il avec $T$ ? » — a deux réponses de signes opposés. À pression constante, chauffer dilate le gaz. À entropie constante (adiabatique réversible), le seul moyen d'augmenter $T$ est de \emph{comprimer} le gaz : $V$ diminue quand $T$ monte. Sans l'indice précisant la variable fixée, l'expression $\partial V/\partial T$ n'a tout simplement pas de sens.

\textbf{Question 4.} En différenciant la loi de Laplace $TV^{\gamma-1} = \text{cste}$ :
\[
V^{\gamma-1}\,\mathrm{d}T + (\gamma-1)TV^{\gamma-2}\,\mathrm{d}V = 0 \implies \left(\frac{\partial V}{\partial T}\right)_{S,N} = -\frac{V}{(\gamma-1)T} = -\frac{Nk_B}{(\gamma-1)P},
\]
qui redonne bien $-\tfrac{3}{2}Nk_B/P$ pour $\gamma = 5/3$. \checkmark\ Le rapport des deux dérivées vaut donc
\[
\frac{\left(\partial V/\partial T\right)_{P,N}}{\left(\partial V/\partial T\right)_{S,N}} = \frac{Nk_B/P}{-Nk_B/[(\gamma-1)P]} = -(\gamma-1) = -\frac{2}{3} \quad (\gamma = 5/3).
\]
\end{solution}
\end{exo}

\begin{exo}[Relations de réciprocité et relation circulaire]{reciprocite-circulaire}
\keywords{relation circulaire, réciprocité, calcul différentiel, variables liées, gaz parfait}

Trois variables $x, y, z$ sont liées par une contrainte $\varphi(x,y,z) = 0$, de sorte que chacune peut s'écrire en fonction des deux autres : $x = x(y,z)$, $y = y(x,z)$, etc.

\begin{enumerate}
\item Écrire $\mathrm{d}x$ en fonction de $\mathrm{d}y$ et $\mathrm{d}z$, puis $\mathrm{d}y$ en fonction de $\mathrm{d}x$ et $\mathrm{d}z$, et substituer la seconde dans la première.
\item En choisissant judicieusement $\mathrm{d}z = 0$ puis $\mathrm{d}x = 0$, démontrer la relation de réciprocité
\[
\left(\frac{\partial x}{\partial y}\right)_z = \frac{1}{\left(\dfrac{\partial y}{\partial x}\right)_z}
\]
puis la relation circulaire
\[
\left(\frac{\partial x}{\partial y}\right)_z\left(\frac{\partial y}{\partial z}\right)_x\left(\frac{\partial z}{\partial x}\right)_y = -1.
\]
\item Vérifier explicitement la relation circulaire pour les variables $(P, V, T)$ d'un gaz parfait.
\item Pourquoi le $-1$ (et non $+1$) est-il contre-intuitif mais correct ?
\end{enumerate}

\begin{indication}
La substitution donne $\mathrm{d}x = \left(\frac{\partial x}{\partial y}\right)_z\left(\frac{\partial y}{\partial x}\right)_z\mathrm{d}x + \left[\left(\frac{\partial x}{\partial y}\right)_z\left(\frac{\partial y}{\partial z}\right)_x + \left(\frac{\partial x}{\partial z}\right)_y\right]\mathrm{d}z$, valable pour $\mathrm{d}x$ et $\mathrm{d}z$ arbitraires : identifier les coefficients.
\end{indication}

\begin{solution}
\textbf{Questions 1--2.}
\[
\mathrm{d}x = \left(\frac{\partial x}{\partial y}\right)_z\mathrm{d}y + \left(\frac{\partial x}{\partial z}\right)_y\mathrm{d}z, \qquad
\mathrm{d}y = \left(\frac{\partial y}{\partial x}\right)_z\mathrm{d}x + \left(\frac{\partial y}{\partial z}\right)_x\mathrm{d}z.
\]
En substituant et en regroupant, l'identité doit valoir pour tout couple $(\mathrm{d}x, \mathrm{d}z)$ indépendant. Le coefficient de $\mathrm{d}x$ donne $1 = \left(\frac{\partial x}{\partial y}\right)_z\left(\frac{\partial y}{\partial x}\right)_z$ : c'est la réciprocité. Le coefficient de $\mathrm{d}z$ donne
\[
\left(\frac{\partial x}{\partial y}\right)_z\left(\frac{\partial y}{\partial z}\right)_x + \left(\frac{\partial x}{\partial z}\right)_y = 0,
\]
qui devient, en divisant par $\left(\frac{\partial x}{\partial z}\right)_y$ et en utilisant la réciprocité, la relation circulaire annoncée.

\textbf{Question 3.} Pour $PV = nRT$ :
\[
\left(\frac{\partial P}{\partial V}\right)_T = -\frac{nRT}{V^2}, \qquad
\left(\frac{\partial V}{\partial T}\right)_P = \frac{nR}{P}, \qquad
\left(\frac{\partial T}{\partial P}\right)_V = \frac{V}{nR}.
\]
Produit : $-\dfrac{nRT}{V^2}\cdot\dfrac{nR}{P}\cdot\dfrac{V}{nR} = -\dfrac{nRT}{PV} = -1$. \checkmark

\textbf{Question 4.} On serait tenté de « simplifier les $\partial$ » comme des fractions, ce qui donnerait $+1$. Mais chaque dérivée est prise à une variable fixée \emph{différente} : ce ne sont pas les mêmes objets, et le produit vaut $-1$. C'est l'illustration la plus frappante du fait que les indices changent la valeur des dérivées partielles.
\end{solution}
\end{exo}

\begin{exo}[Construire l'enthalpie et l'enthalpie libre]{legendre-enthalpie}
\keywords{transformée de Legendre, enthalpie, enthalpie libre, fonction de Gibbs, variables naturelles, relations de Maxwell}

On part du premier principe pour un système ouvert en transformation réversible : $\mathrm{d}U = T\,\mathrm{d}S - P\,\mathrm{d}V + \mu\,\mathrm{d}N$.

\begin{enumerate}
\item On cherche une fonction d'état $H$ dont les variables naturelles soient $(S, P, N)$. Montrer que $H = U + PV$ convient, écrire $\mathrm{d}H$, et en déduire les trois équations d'état associées.
\item Écrire la relation de Maxwell issue des dérivées croisées de $H$ par rapport à $S$ et $P$.
\item Construire de même la fonction $G$ de variables naturelles $(T, P, N)$ et écrire $\mathrm{d}G$. Pourquoi $G$ est-elle le potentiel adapté à la chimie ?
\item Pour un gaz parfait monoatomique, exprimer $H$ en fonction de $T$ et $N$ uniquement. En déduire $\left(\partial H/\partial T\right)_{P,N}$ et faire le lien avec la chaleur reçue lors d'une transformation isobare.
\end{enumerate}

\begin{indication}
Transformée de Legendre : pour échanger une variable extensive contre sa conjuguée intensive, ajouter ou retrancher le produit des deux. Pour la question 4, $H = U + PV = \frac{3}{2}Nk_BT + Nk_BT$.
\end{indication}

\begin{solution}
\textbf{Question 1.} $\mathrm{d}H = \mathrm{d}U + P\,\mathrm{d}V + V\,\mathrm{d}P = T\,\mathrm{d}S + V\,\mathrm{d}P + \mu\,\mathrm{d}N$ : les variables naturelles sont bien $(S,P,N)$, et par identification
\[
T = \left(\frac{\partial H}{\partial S}\right)_{P,N}, \qquad V = \left(\frac{\partial H}{\partial P}\right)_{S,N}, \qquad \mu = \left(\frac{\partial H}{\partial N}\right)_{S,P}.
\]

\textbf{Question 2.} Schwarz sur les deux premières : $\left(\dfrac{\partial T}{\partial P}\right)_{S,N} = \left(\dfrac{\partial V}{\partial S}\right)_{P,N}$.

\textbf{Question 3.} Il faut de plus échanger $S$ contre $T$ : $G = H - TS = U + PV - TS$, d'où
\[
\mathrm{d}G = -S\,\mathrm{d}T + V\,\mathrm{d}P + \mu\,\mathrm{d}N.
\]
À $T$ et $P$ constants (conditions usuelles d'une paillasse de chimie), $\mathrm{d}G = \mu\,\mathrm{d}N$ : la variation de $G$ ne contient que la chimie, c'est-à-dire les variations des nombres de particules pondérées par les potentiels chimiques.

\textbf{Question 4.} $H = U + PV = \tfrac{3}{2}Nk_BT + Nk_BT = \tfrac{5}{2}Nk_BT$ : comme $U$, l'enthalpie du gaz parfait ne dépend que de $T$ (et de $N$). Donc $\left(\partial H/\partial T\right)_{P,N} = \tfrac{5}{2}Nk_B = C_P$. Et pour une transformation isobare d'un système fermé, $\mathrm{d}H = T\,\mathrm{d}S + V\,\mathrm{d}P = \delta Q$ : la variation d'enthalpie \emph{est} la chaleur reçue — on retrouve l'expérience de la résistance sous piston de la leçon 5, où $Q = \tfrac{5}{2}R\,\Delta T$ par mole.
\end{solution}
\end{exo}

\begin{exo}[Relation de Maxwell à partir de l'énergie libre de Helmholtz]{maxwell-helmholtz}
\keywords{relation de Maxwell, énergie libre, potentiel thermodynamique, dérivées secondes, gaz parfait}

L'énergie libre de Helmholtz $F$ est définie par $F = U - TS$. Sa différentielle est :
$$\mathrm{d}F = -S\,\mathrm{d}T - P\,\mathrm{d}V.$$

\begin{enumerate}
  \item Identifier les dérivées partielles de $F$ par rapport à $T$ et $V$.
  \item En utilisant la symétrie des dérivées secondes croisées, établir la relation de Maxwell associée à $F$.
  \item Vérifier cette relation pour un gaz parfait ($P = nRT/V$).
  \item Interpréter physiquement la grandeur $\left(\partial P/\partial T\right)_V$.
\end{enumerate}

\begin{indication}
Si $\mathrm{d}F = A\,\mathrm{d}T + B\,\mathrm{d}V$, alors $\left(\partial A/\partial V\right)_T = \left(\partial B/\partial T\right)_V$ (théorème de Schwarz).
\end{indication}

\begin{solution}
\textbf{1.} Par identification : $\left(\partial F/\partial T\right)_V = -S$ et $\left(\partial F/\partial V\right)_T = -P$.

\textbf{2.} Schwarz appliqué à $F$ :
$$\frac{\partial^2 F}{\partial V\,\partial T} = \frac{\partial^2 F}{\partial T\,\partial V} \implies \left(\frac{\partial S}{\partial V}\right)_T = \left(\frac{\partial P}{\partial T}\right)_V.$$

\textbf{3.} Pour un gaz parfait : $\left(\partial P/\partial T\right)_V = nR/V$. Et $S = nC_v\ln T + nR\ln V + \text{cste}$, donc $\left(\partial S/\partial V\right)_T = nR/V$. Les deux membres sont bien égaux. \checkmark

\textbf{4.} $\left(\partial P/\partial T\right)_V$ mesure comment la pression augmente avec la température à volume constant — c'est directement lié au coefficient d'expansion thermique et à la compressibilité du fluide.
\end{solution}

\end{exo}

\begin{exo}[Du potentiel de Helmholtz à la pression de rayonnement]{helmholtz-photons-voile-solaire}
\keywords{gaz de photons, énergie libre de Helmholtz, rayonnement du corps noir, pression de rayonnement, voile solaire}

Un rayonnement électromagnétique à l'équilibre dans une cavité constitue un gaz de photons. On admet directement son énergie libre de Helmholtz,
\[
F(T,V)=-\frac{a_{\rm r}}{3}VT^4,
\qquad
a_{\rm r}=7{,}57\times10^{-16}\,\mathrm{J\,m^{-3}\,K^{-4}}.
\]
Le nombre de photons n'étant pas conservé, $F$ ne dépend pas d'une variable $N$. On rappelle
\[
\mathrm dF=-S\,\mathrm dT-P\,\mathrm dV.
\]

\begin{enumerate}
\item Déduire de $F(T,V)$ l'entropie $S$, la pression $P$, puis l'énergie interne $U=F+TS$. Retrouver $P=U/(3V)$.
\item Vérifier la relation d'Euler $U=TS-PV$ et montrer que $F=-PV$. Pourquoi aucun terme $\mu N$ n'apparaît-il ?
\item Estimer la densité d'énergie et la pression du rayonnement d'une cavité à $T=300\,\mathrm K$. Commenter l'ordre de grandeur.
\item Une voile solaire parfaitement réfléchissante de surface $A=100\,\mathrm{m^2}$ et de masse totale $m=10\,\mathrm{kg}$ reçoit normalement, près de la Terre, un faisceau d'intensité $I=1{,}36\times10^3\,\mathrm{W\,m^{-2}}$. Pour un faisceau dirigé, on admet que la pression vaut $I/c$ s'il est absorbé et $2I/c$ s'il est réfléchi, avec $c=3{,}00\times10^8\,\mathrm{m\,s^{-1}}$. Calculer la force, l'accélération et le gain de vitesse en une journée si l'accélération reste constante. Pourquoi n'utilise-t-on pas ici le facteur $1/3$ de la cavité ?
\end{enumerate}

\begin{indication}
Utiliser $S=-(\partial F/\partial T)_V$, $P=-(\partial F/\partial V)_T$ et $U=F+TS$. Pour la voile, la force est la pression multipliée par la surface et une journée vaut $86400\,\mathrm s$.
\end{indication}

\begin{solution}
\textbf{Question 1.} Les variables naturelles de $F$ donnent
\[
S=-\left(\frac{\partial F}{\partial T}\right)_V
=\frac{4a_{\rm r}}{3}VT^3,
\qquad
P=-\left(\frac{\partial F}{\partial V}\right)_T
=\frac{a_{\rm r}}{3}T^4.
\]
Par conséquent,
\[
U=F+TS=-\frac{a_{\rm r}}{3}VT^4
+\frac{4a_{\rm r}}{3}VT^4
=\boxed{a_{\rm r}VT^4},
\]
et donc $P=a_{\rm r}T^4/3=U/(3V)$.

\textbf{Question 2.}
\[
TS-PV=\frac{4a_{\rm r}}{3}VT^4-\frac{a_{\rm r}}{3}VT^4
=a_{\rm r}VT^4=U.
\]
De plus, $F=U-TS=-PV$. Le nombre de photons s'ajuste librement à l'équilibre : il n'est pas une variable extensive conservée. Il n'y a donc pas de terme $\mu\,\mathrm dN$ indépendant ; on dit de manière équivalente que le potentiel chimique des photons est nul.

\textbf{Question 3.} La densité d'énergie vaut
\[
u=\frac UV=a_{\rm r}T^4
=7{,}57\times10^{-16}\times300^4
\simeq6{,}1\times10^{-6}\,\mathrm{J\,m^{-3}}.
\]
Ainsi,
\[
P=\frac u3\simeq2{,}0\times10^{-6}\,\mathrm{Pa}.
\]
À température ambiante, la pression du rayonnement thermique est minuscule devant la pression atmosphérique, de l'ordre de $10^5\,\mathrm{Pa}$. Elle devient toutefois importante à très haute température, notamment à l'intérieur des étoiles.

\textbf{Question 4.} Pour une réflexion parfaite,
\[
P_{\rm voile}=\frac{2I}{c}
\simeq9{,}1\times10^{-6}\,\mathrm{Pa}.
\]
La force et l'accélération valent
\[
F=P_{\rm voile}A\simeq9{,}1\times10^{-4}\,\mathrm N,
\qquad
a=\frac Fm\simeq9{,}1\times10^{-5}\,\mathrm{m\,s^{-2}}.
\]
En une journée,
\[
\Delta v=a\,\Delta t
\simeq9{,}1\times10^{-5}\times86400
\simeq7{,}8\,\mathrm{m\,s^{-1}}.
\]
La force instantanée est faible, mais elle agit sans carburant et son effet s'accumule. Le facteur $1/3$ caractérise un rayonnement isotrope en équilibre dans une cavité : les photons arrivent de toutes les directions. La voile reçoit au contraire un faisceau presque unidirectionnel ; le transfert de quantité de mouvement normal vaut $I/c$, doublé par la réflexion.
\end{solution}
\end{exo}

\begin{exo}[Le piège du potentiel chimique du gaz parfait]{potentiel-chimique-gp}
\keywords{potentiel chimique, relation d'Euler, enthalpie libre, dérivée partielle, variables naturelles, erreur classique}

Un étudiant raisonne ainsi : « puisque $\mu = \left(\partial U/\partial N\right)_{S,V}$ et que $U = \frac{3}{2}Nk_BT$ pour le gaz parfait monoatomique, alors $\mu = \frac{3}{2}k_BT$. »

\begin{enumerate}
\item À partir de la relation d'Euler $U = TS - PV + \mu N$, montrer que $G = U - TS + PV = \mu N$ : pour un corps pur, $\mu$ est l'enthalpie libre par particule.
\item En déduire l'expression correcte de $\mu$ pour le gaz parfait monoatomique, en fonction de $T$, $S$, $N$ et $k_B$. Comparer au résultat de l'exercice \textit{Sackur-Tétrode}.
\item Où est l'erreur dans le raisonnement de l'étudiant ? Écrire le terme qu'il a oublié.
\item Le résultat correct peut-il coïncider avec $\frac{3}{2}k_BT$ ? Commenter le signe de $\mu$ pour un gaz parfait usuel.
\end{enumerate}

\begin{indication}
Dans la dérivée $\left(\partial U/\partial N\right)_{S,V}$, ce sont $S$ et $V$ qui sont fixés — pas $T$. Or $T = T(S,V,N)$ varie quand on ajoute des particules à $S,V$ fixés.
\end{indication}

\begin{solution}
\textbf{Question 1.} L'extensivité (théorème d'Euler) donne $U = TS - PV + \mu N$, donc $G = U - TS + PV = \mu N$ directement : $\mu = G/N = g$, l'enthalpie libre par particule.

\textbf{Question 2.} Pour le gaz parfait monoatomique, $U = \tfrac{3}{2}Nk_BT$ et $PV = Nk_BT$ :
\[
\mu = \frac{U - TS + PV}{N} = \frac{5}{2}k_BT - T\,\frac{S}{N},
\]
exactement le résultat obtenu en dérivant Sackur-Tétrode. \checkmark

\textbf{Question 3.} En écrivant $\partial(\tfrac{3}{2}Nk_BT)/\partial N = \tfrac{3}{2}k_BT$, l'étudiant suppose $T$ indépendante de $N$. Mais dans les variables naturelles $(S,V,N)$ de $U$, la température est une fonction $T(S,V,N)$ : ajouter des particules à entropie et volume fixés change la température. Le calcul correct s'écrit
\[
\mu = \left(\frac{\partial U}{\partial N}\right)_{S,V} = \frac{3}{2}k_BT + \frac{3}{2}Nk_B\left(\frac{\partial T}{\partial N}\right)_{S,V},
\]
et le second terme n'est pas nul.

\textbf{Question 4.} Il faudrait $S/N = k_B$, ce qui n'est jamais le cas pour un gaz usuel : l'entropie par particule vaut typiquement plusieurs dizaines de $k_B$ (via Sackur-Tétrode). Du coup $\mu = \tfrac{5}{2}k_BT - T S/N < 0$ : le potentiel chimique du gaz parfait est négatif — ajouter une particule à $U$ et $V$ fixés \emph{augmente} l'entropie, donc l'équilibre accepte volontiers de nouvelles particules.
\end{solution}
\end{exo}

\begin{exo}[Toute la thermodynamique à partir de F]{energie-libre-berthelot}
\keywords{énergie libre, gaz réel, équation d'état, entropie, capacité thermique, relation de Maxwell, gaz de Berthelot}

L'énergie libre molaire d'un gaz réel (gaz de Berthelot) s'écrit
\[
F(T,V) = -RT\ln\!\left(\frac{V-b}{b}\right) - \varphi(T) - \frac{a}{VT},
\]
où $a$ et $b$ sont des constantes positives et $\varphi(T)$ une fonction de la température seule.

\begin{enumerate}
\item Justifier la forme différentielle $\mathrm{d}F = -S\,\mathrm{d}T - P\,\mathrm{d}V$ (nombre de moles fixé).
\item En déduire l'équation d'état $P(T,V)$ de ce gaz.
\item Déterminer $S(T,V)$, puis montrer que $U = T\varphi'(T) - \varphi(T) - \dfrac{2a}{VT}$.
\item Calculer $C_V = \left(\partial U/\partial T\right)_V$. Sachant que $C_V$ tend vers une constante $C$ quand $V \to \infty$, déterminer la forme générale de $\varphi(T)$. On donne $\int_0^T\ln\theta\,\mathrm{d}\theta = T(\ln T - 1)$.
\item Vérifier la relation de Maxwell $\left(\partial S/\partial V\right)_T = \left(\partial P/\partial T\right)_V$ sur les expressions obtenues.
\item Que devient ce gaz dans la limite $a \to 0$ ? Interpréter les rôles respectifs de $a$ et $b$.
\end{enumerate}

\begin{indication}
$F = U - TS$ et le premier principe donnent $\mathrm{d}F = -S\,\mathrm{d}T - P\,\mathrm{d}V$, d'où $P = -\left(\partial F/\partial V\right)_T$ et $S = -\left(\partial F/\partial T\right)_V$, puis $U = F + TS$.
\end{indication}

\begin{solution}
\textbf{Question 1.} $\mathrm{d}F = \mathrm{d}U - T\,\mathrm{d}S - S\,\mathrm{d}T$, et pour une transformation réversible d'un système fermé $\mathrm{d}U = T\,\mathrm{d}S - P\,\mathrm{d}V$, donc $\mathrm{d}F = -S\,\mathrm{d}T - P\,\mathrm{d}V$.

\textbf{Question 2.}
\[
P = -\left(\frac{\partial F}{\partial V}\right)_T = \frac{RT}{V-b} - \frac{a}{TV^2}.
\]
C'est une variante de van der Waals où le terme attractif décroît avec la température.

\textbf{Question 3.}
\[
S = -\left(\frac{\partial F}{\partial T}\right)_V = R\ln\!\left(\frac{V-b}{b}\right) + \varphi'(T) - \frac{a}{VT^2},
\]
puis $U = F + TS = T\varphi'(T) - \varphi(T) - \dfrac{2a}{VT}$ (les logarithmes se compensent).

\textbf{Question 4.} $C_V = \left(\partial U/\partial T\right)_V = T\varphi''(T) + \dfrac{2a}{VT^2}$. Quand $V\to\infty$, $C_V \to T\varphi''(T) = C$, donc $\varphi''(T) = C/T$, soit $\varphi'(T) = C\ln T + K_1$ et, avec l'intégrale donnée,
\[
\varphi(T) = CT\ln T + (K_1 - C)\,T + K_2,
\]
que l'on peut réécrire $\varphi(T) = CT\ln T + K_1'T + K_2$ en renommant la constante.

\textbf{Question 5.} $\left(\dfrac{\partial S}{\partial V}\right)_T = \dfrac{R}{V-b} + \dfrac{a}{V^2T^2}$ et $\left(\dfrac{\partial P}{\partial T}\right)_V = \dfrac{R}{V-b} + \dfrac{a}{TV^2}\cdot\dfrac{1}{T} = \dfrac{R}{V-b} + \dfrac{a}{V^2T^2}$. Égalité. \checkmark

\textbf{Question 6.} Pour $a \to 0$ : $P = RT/(V-b)$, $U = T\varphi' - \varphi$ ne dépend plus que de $T$ : on retrouve un gaz parfait à covolume près. Le paramètre $b$ (covolume) traduit le volume propre des molécules, tandis que $a$ traduit leur attraction mutuelle — c'est lui qui fait dépendre $U$ du volume, signature d'un gaz réel.
\end{solution}
\end{exo}
